% DISCUSSION, LIMITATIONS AND THREATS TO VALIDITY. Numbers ONLY via \R... macros.
\subsection{WHAT THE RESULTS MEAN}
The evidence separates two uses of measured context value. As a \emph{selector} the MCV compiler was no better than
embedding relevance, and worse when its block-level ranking learned the wrong correlate. As a \emph{type-level
budgeter and diagnostic} it was useful. It showed that worked examples were dead weight in every family, that a stale note
was redundant, and that tool schemas lose most of their value when shortened. Acting on that, it cut prompt tokens
substantially on tool use without losing success. The profile also stopped spending once no remaining block had
positive value, which a fill-the-budget selector does not do. These findings suggest a hybrid in which relevance ranks
blocks within a type and the profile decides which types to include and when to stop. We did not test that hybrid, and
we state it only as the design the data point to.

The negative result is informative about leave-one-block-out supervision. When many blocks are individually
dispensable (one session of \RNumSessions, one document of \RNumDocs), single deletions rarely change the outcome. A
regression on those deltas then learns proxies such as recency and length. Group or pairwise deletions, or using
relevance directly as the within-type ranker, are natural fixes.

\subsection{PRACTICAL GUIDANCE}
For practitioners the results suggest a cheap procedure. Profile each (model, task family) on a small development set
(ten to twenty instances sufficed here to rank types) using type removals and the length-neutral control. Remove types
whose value interval sits at zero, such as the worked examples here, from the default prompt. Keep a strong relevance
ranker for choosing blocks within a type. Stop adding context when the remaining types have no measured value rather than
filling the budget. Re-profile when the model or the task family changes; the 4B-to-8B agreement observed here is
encouraging but limited to one model family. Interaction terms can be omitted unless the profile shows a large
interaction between two types that both carry value.

\subsection{LIMITATIONS AND THREATS TO VALIDITY}
\textbf{Construct validity.} Each instance is a single agent decision, not a multi-step trajectory; context value in
long-horizon agents may compound across steps. Success is a deterministic check (answer match, AST-level call match),
which misses partially correct outputs.
\textbf{Internal validity.} B4 is our re-implementation of a PACMS-style objective, not the authors' code, and B3 uses one
LLMLingua-2 checkpoint; both could be tuned further. Rendered prompts may exceed the budget by at most \RMaxOvershoot
tokens because tokenization is not additive across block boundaries. Greedy decoding on Apple-silicon inference is not
guaranteed to be bitwise deterministic; caching makes every reported response reproducible from the logs, but a fresh
rerun may differ slightly.
\textbf{Statistical conclusion validity.} The pre-registered power analysis put power for a \RSesoi{} AUBC effect on facts at
about \RPowerFacts with \RNMain instances, so the null result on facts is weak evidence of no effect. The weaker baselines used
\RNWeak instances; each scored below B2 in every family (Table~\ref{tab:main}), so B2 is the best baseline throughout. Transfer used
\RNTransfer instances per family and two budgets.
\textbf{External validity.} Two models of one family (Qwen3, 4B and 8B) were tested; a third model from another family
was planned and dropped for compute, so transfer across model families is untested. One task family (history) is
synthetic, and HotpotQA and BFCL may overlap with pretraining data. Profiles are per (model, task family) and must be
re-measured when either changes, at the profiling cost reported in Section~\ref{sec:eval}.
